% ===== VARIANT: QUANTITATIVE RESEARCH ===============================
% Angle: methodological originality, proofs, calibration theory,
% stochastic control, filtering. Research output leads; applied credit
% work compressed.
% ===================================================================
% =====================================================================
%  Shared preamble for all CV variants.
%  Each variant does:   % =====================================================================
%  Shared preamble for all CV variants.
%  Each variant does:   % =====================================================================
%  Shared preamble for all CV variants.
%  Each variant does:   \input{cv-preamble.tex}
%  Do not put content here — only packages, formatting and macros.
% =====================================================================
\documentclass[letterpaper,11pt]{article}

\usepackage{latexsym}
\usepackage[empty]{fullpage}
\usepackage{titlesec}
\usepackage{marvosym}
\usepackage[usenames,dvipsnames]{color}
\usepackage{verbatim}
\usepackage{enumitem}
\usepackage[hidelinks]{hyperref}
\usepackage{fancyhdr}
\usepackage[english]{babel}
\usepackage{tabularx}
\usepackage{fontawesome5}
\usepackage{multicol}
\usepackage{graphicx}
\setlength{\multicolsep}{-3.0pt}
\setlength{\columnsep}{-1pt}
\input{glyphtounicode}
\usepackage{paracol}

\pagestyle{fancy}
\fancyhf{}
\fancyfoot{}
\renewcommand{\headrulewidth}{0pt}
\renewcommand{\footrulewidth}{0pt}

\addtolength{\oddsidemargin}{-0.5in}
\addtolength{\evensidemargin}{-0.5in}
\addtolength{\textwidth}{0.9in}
\addtolength{\topmargin}{-.3in}
\addtolength{\textheight}{0.9in}
\hypersetup{
    colorlinks=true,
    linkcolor=blue,
    urlcolor=blue,
    filecolor=blue,
    pdfpagemode=FullScreen
}
\urlstyle{same}

\raggedbottom
\raggedright
\setlength{\tabcolsep}{0in}

\titleformat{\section}{
  \vspace{-4pt}\scshape\raggedright\large\bfseries
}{}{0em}{}[\color{black}\titlerule \vspace{-5pt}]

\pdfgentounicode=1

\newcommand{\resumeItem}[1]{
  \item\small{{#1 \vspace{-2pt}}}
}

\newcommand{\resumeSubheading}[4]{
  \vspace{-2pt}\item
    \begin{tabular*}{1.0\textwidth}[t]{l@{\extracolsep{\fill}}r}
      \textbf{#1} & \textbf{\small #2} \\
      \textit{\small#3} & \textit{\small #4} \\
    \end{tabular*}\vspace{-7pt}
}

\newcommand{\resumeProjectHeading}[2]{
    \item
    \begin{tabular*}{0.97\textwidth}{l@{\extracolsep{\fill}}r}
      \small#1 & \textbf{\small #2}\\
    \end{tabular*}\vspace{-7pt}
}

\renewcommand\labelitemi{$\vcenter{\hbox{\tiny$\bullet$}}$}
\renewcommand\labelitemii{$\vcenter{\hbox{\tiny$\bullet$}}$}

\newcommand{\resumeSubHeadingListStart}{\begin{itemize}[leftmargin=0.0in, label={}]}
\newcommand{\resumeSubHeadingListEnd}{\end{itemize}}
\newcommand{\resumeItemListStart}{\begin{itemize}}
\newcommand{\resumeItemListEnd}{\end{itemize}\vspace{-5pt}}

% ---------------------------------------------------------------------
% \cvheader{positioning line}{focus line}
% ---------------------------------------------------------------------
\newcommand{\cvheader}[2]{%
\begin{center}
    {\Huge \scshape Franck Deba} \\[0.4em]
    {\small \textit{#1}} \\[0.25em]
    {\small #2}
\end{center}
\vspace{0.5pt}
\small
\raisebox{-0.1\height}{\faPhone}~+33 (0)6 05 77 47 51 \hspace{1em}
\raisebox{-0.2\height}{\faEnvelope}~\href{mailto:debafranck@gmail.com}{\underline{debafranck@gmail.com}} \hspace{1em}
\raisebox{-0.2\height}{\faLinkedin}~\href{https://linkedin.com/in/deba-franck}{\underline{deba-franck}} \hspace{1em}
\raisebox{-0.2\height}{\faGithub}~\href{https://github.com/usenameFD}{\underline{usenameFD}} \hspace{1em}
\raisebox{-0.2\height}{\faGlobe}~\href{https://usenamefd.github.io/}{\underline{usenamefd.github.io}}
\vspace{-8pt}
}

%  Do not put content here — only packages, formatting and macros.
% =====================================================================
\documentclass[letterpaper,11pt]{article}

\usepackage{latexsym}
\usepackage[empty]{fullpage}
\usepackage{titlesec}
\usepackage{marvosym}
\usepackage[usenames,dvipsnames]{color}
\usepackage{verbatim}
\usepackage{enumitem}
\usepackage[hidelinks]{hyperref}
\usepackage{fancyhdr}
\usepackage[english]{babel}
\usepackage{tabularx}
\usepackage{fontawesome5}
\usepackage{multicol}
\usepackage{graphicx}
\setlength{\multicolsep}{-3.0pt}
\setlength{\columnsep}{-1pt}
\input{glyphtounicode}
\usepackage{paracol}

\pagestyle{fancy}
\fancyhf{}
\fancyfoot{}
\renewcommand{\headrulewidth}{0pt}
\renewcommand{\footrulewidth}{0pt}

\addtolength{\oddsidemargin}{-0.5in}
\addtolength{\evensidemargin}{-0.5in}
\addtolength{\textwidth}{0.9in}
\addtolength{\topmargin}{-.3in}
\addtolength{\textheight}{0.9in}
\hypersetup{
    colorlinks=true,
    linkcolor=blue,
    urlcolor=blue,
    filecolor=blue,
    pdfpagemode=FullScreen
}
\urlstyle{same}

\raggedbottom
\raggedright
\setlength{\tabcolsep}{0in}

\titleformat{\section}{
  \vspace{-4pt}\scshape\raggedright\large\bfseries
}{}{0em}{}[\color{black}\titlerule \vspace{-5pt}]

\pdfgentounicode=1

\newcommand{\resumeItem}[1]{
  \item\small{{#1 \vspace{-2pt}}}
}

\newcommand{\resumeSubheading}[4]{
  \vspace{-2pt}\item
    \begin{tabular*}{1.0\textwidth}[t]{l@{\extracolsep{\fill}}r}
      \textbf{#1} & \textbf{\small #2} \\
      \textit{\small#3} & \textit{\small #4} \\
    \end{tabular*}\vspace{-7pt}
}

\newcommand{\resumeProjectHeading}[2]{
    \item
    \begin{tabular*}{0.97\textwidth}{l@{\extracolsep{\fill}}r}
      \small#1 & \textbf{\small #2}\\
    \end{tabular*}\vspace{-7pt}
}

\renewcommand\labelitemi{$\vcenter{\hbox{\tiny$\bullet$}}$}
\renewcommand\labelitemii{$\vcenter{\hbox{\tiny$\bullet$}}$}

\newcommand{\resumeSubHeadingListStart}{\begin{itemize}[leftmargin=0.0in, label={}]}
\newcommand{\resumeSubHeadingListEnd}{\end{itemize}}
\newcommand{\resumeItemListStart}{\begin{itemize}}
\newcommand{\resumeItemListEnd}{\end{itemize}\vspace{-5pt}}

% ---------------------------------------------------------------------
% \cvheader{positioning line}{focus line}
% ---------------------------------------------------------------------
\newcommand{\cvheader}[2]{%
\begin{center}
    {\Huge \scshape Franck Deba} \\[0.4em]
    {\small \textit{#1}} \\[0.25em]
    {\small #2}
\end{center}
\vspace{0.5pt}
\small
\raisebox{-0.1\height}{\faPhone}~+33 (0)6 05 77 47 51 \hspace{1em}
\raisebox{-0.2\height}{\faEnvelope}~\href{mailto:debafranck@gmail.com}{\underline{debafranck@gmail.com}} \hspace{1em}
\raisebox{-0.2\height}{\faLinkedin}~\href{https://linkedin.com/in/deba-franck}{\underline{deba-franck}} \hspace{1em}
\raisebox{-0.2\height}{\faGithub}~\href{https://github.com/usenameFD}{\underline{usenameFD}} \hspace{1em}
\raisebox{-0.2\height}{\faGlobe}~\href{https://usenamefd.github.io/}{\underline{usenamefd.github.io}}
\vspace{-8pt}
}

%  Do not put content here — only packages, formatting and macros.
% =====================================================================
\documentclass[letterpaper,11pt]{article}

\usepackage{latexsym}
\usepackage[empty]{fullpage}
\usepackage{titlesec}
\usepackage{marvosym}
\usepackage[usenames,dvipsnames]{color}
\usepackage{verbatim}
\usepackage{enumitem}
\usepackage[hidelinks]{hyperref}
\usepackage{fancyhdr}
\usepackage[english]{babel}
\usepackage{tabularx}
\usepackage{fontawesome5}
\usepackage{multicol}
\usepackage{graphicx}
\setlength{\multicolsep}{-3.0pt}
\setlength{\columnsep}{-1pt}
\input{glyphtounicode}
\usepackage{paracol}

\pagestyle{fancy}
\fancyhf{}
\fancyfoot{}
\renewcommand{\headrulewidth}{0pt}
\renewcommand{\footrulewidth}{0pt}

\addtolength{\oddsidemargin}{-0.5in}
\addtolength{\evensidemargin}{-0.5in}
\addtolength{\textwidth}{0.9in}
\addtolength{\topmargin}{-.3in}
\addtolength{\textheight}{0.9in}
\hypersetup{
    colorlinks=true,
    linkcolor=blue,
    urlcolor=blue,
    filecolor=blue,
    pdfpagemode=FullScreen
}
\urlstyle{same}

\raggedbottom
\raggedright
\setlength{\tabcolsep}{0in}

\titleformat{\section}{
  \vspace{-4pt}\scshape\raggedright\large\bfseries
}{}{0em}{}[\color{black}\titlerule \vspace{-5pt}]

\pdfgentounicode=1

\newcommand{\resumeItem}[1]{
  \item\small{{#1 \vspace{-2pt}}}
}

\newcommand{\resumeSubheading}[4]{
  \vspace{-2pt}\item
    \begin{tabular*}{1.0\textwidth}[t]{l@{\extracolsep{\fill}}r}
      \textbf{#1} & \textbf{\small #2} \\
      \textit{\small#3} & \textit{\small #4} \\
    \end{tabular*}\vspace{-7pt}
}

\newcommand{\resumeProjectHeading}[2]{
    \item
    \begin{tabular*}{0.97\textwidth}{l@{\extracolsep{\fill}}r}
      \small#1 & \textbf{\small #2}\\
    \end{tabular*}\vspace{-7pt}
}

\renewcommand\labelitemi{$\vcenter{\hbox{\tiny$\bullet$}}$}
\renewcommand\labelitemii{$\vcenter{\hbox{\tiny$\bullet$}}$}

\newcommand{\resumeSubHeadingListStart}{\begin{itemize}[leftmargin=0.0in, label={}]}
\newcommand{\resumeSubHeadingListEnd}{\end{itemize}}
\newcommand{\resumeItemListStart}{\begin{itemize}}
\newcommand{\resumeItemListEnd}{\end{itemize}\vspace{-5pt}}

% ---------------------------------------------------------------------
% \cvheader{positioning line}{focus line}
% ---------------------------------------------------------------------
\newcommand{\cvheader}[2]{%
\begin{center}
    {\Huge \scshape Franck Deba} \\[0.4em]
    {\small \textit{#1}} \\[0.25em]
    {\small #2}
\end{center}
\vspace{0.5pt}
\small
\raisebox{-0.1\height}{\faPhone}~+33 (0)6 05 77 47 51 \hspace{1em}
\raisebox{-0.2\height}{\faEnvelope}~\href{mailto:debafranck@gmail.com}{\underline{debafranck@gmail.com}} \hspace{1em}
\raisebox{-0.2\height}{\faLinkedin}~\href{https://linkedin.com/in/deba-franck}{\underline{deba-franck}} \hspace{1em}
\raisebox{-0.2\height}{\faGithub}~\href{https://github.com/usenameFD}{\underline{usenameFD}} \hspace{1em}
\raisebox{-0.2\height}{\faGlobe}~\href{https://usenamefd.github.io/}{\underline{usenamefd.github.io}}
\vspace{-8pt}
}


\begin{document}

\cvheader{Quantitative Researcher}{Stochastic control $\cdot$ Factor model calibration $\cdot$ Filtering $\cdot$ Numerical methods for PDEs $\cdot$ Risk modelling}

%-----------RESEARCH-----------
\section{Research}
\resumeSubHeadingListStart

  \resumeProjectHeading
    {\href{https://usenamefd.github.io/projects/drc-frtb-cac40.html}{\textbf{Factor-Model Calibration for the Default Risk Charge}} $|$ \emph{M2 thesis, 47 pp.}} {2026}
    \resumeItemListStart
      \resumeItem{Compared three calibration criteria for the latent multi-factor model underlying the DRC: iterative PCA, Frobenius-norm minimisation and maximum likelihood --- the latter being, to my knowledge, a new contribution to DRC modelling}
      \resumeItem{Showed iterative PCA and Frobenius converge to an identical fit, while the likelihood criterion --- equivalent to minimising a Kullback--Leibler divergence --- penalises correlation under-estimation asymmetrically, concentrating its effect on the most correlated issuer block}
      \resumeItem{Frobenius and MLE share one spectral projected-gradient solver, so measured gaps isolate the criterion rather than the implementation}
      \resumeItem{Established that maximum likelihood is insensitive to the choice of empirical-correlation estimator (raw vs.\ shrunk), unlike the two matching methods}
      \resumeItem{Monte Carlo DRC with 200{,}000 scenarios and sectioning confidence intervals; Marchenko--Pastur independently selects two factors on all calibration windows}
      \resumeItem{Systematic stationarity (ADF, KPSS) and normality (Shapiro--Wilk, Kolmogorov--Smirnov) testing of factors and residuals}
    \resumeItemListEnd

  \resumeProjectHeading
    {\href{https://github.com/usenameFD/Numerical_methods/blob/main/zermeloNotebook.ipynb}{\textbf{Numerical Resolution of an HJBI PDE --- Zermelo Stochastic Control Problem}} $|$ \emph{Python, PyTorch}} {2026}
    \resumeItemListStart
      \resumeItem{Solved a Hamilton--Jacobi--Bellman--Isaacs PDE arising from a stochastic optimal control problem}
      \resumeItem{Finite-difference scheme combined with a semi-smooth Newton method; measured numerical order of convergence against a known closed-form solution before applying it to a case with no explicit solution}
      \resumeItem{Physics-Informed Neural Network solver using automatic differentiation (Esteve-Yag\"ue, 2025); compared three loss formulations and studied hyperparameter sensitivity}
      \resumeItem{Assessed accuracy, convergence and robustness --- PINNs promising in higher dimensions, finite differences retaining a provable convergence order}
    \resumeItemListEnd

  \resumeProjectHeading
    {\href{https://github.com/usenameFD/Conjecture-de-Collatz}{\textbf{An Analytic Reformulation of the Collatz Conjecture}}} {2026}
    \resumeItemListStart
      \resumeItem{Reformulated the conjecture as a uniqueness problem for the stationary measure of a Markov chain}
      \resumeItem{Numerical verification of the characterisation of admissible stationary measures}
    \resumeItemListEnd

  \resumeProjectHeading
    {\href{https://usenamefd.github.io/projects/methodes-numeriques-edp.html}{\textbf{Optimal Stopping and Convergence to the Brownian Bridge}} $|$ \emph{Python}} {2026}
    \resumeItemListStart
      \resumeItem{Modelled a card-game brainteaser (Crack, 2014) as an optimal stopping problem solved by backward induction}
      \resumeItem{Computed the discrete stopping boundary and proved convergence to the continuous boundary; asymptotic analysis of the normalised value following Ekstr\"om (2020)}
    \resumeItemListEnd

  \resumeProjectHeading
    {\href{https://usenamefd.github.io/projects/calibration-heston-filtrage.html}{\textbf{Stochastic Volatility Calibration by Particle Filtering}} $|$ \emph{Python, R}} {2025}
    \resumeItemListStart
      \resumeItem{Heston option pricing by Fourier inversion; variance treated as a hidden state in a state-space formulation using the exact non-central chi-square transition law}
      \resumeItem{Implemented bootstrap and auxiliary particle filters; simulation study comparing filter accuracy by option moneyness (ATM, ITM, mixed)}
      \resumeItem{Companion work on the Taylor log-SV model: linearisation with a Kalman filter vs.\ a bootstrap particle filter, quantifying the cost of the Gaussian approximation of a log-$\chi^2_1$ noise}
    \resumeItemListEnd

  \resumeProjectHeading
    {\href{https://usenamefd.github.io/projects/creditvar-copules.html}{\textbf{Dependence Structures and Credit Portfolio Risk}} $|$ \emph{Python}} {2025}
    \resumeItemListStart
      \resumeItem{Separated marginal and dependence estimation (Sklar); skewed-Student margins by maximum likelihood, elliptical and Archimedean copulas compared by kernel-density diagnostics}
      \resumeItem{Empirical evidence of upper-tail dependence incompatible with the asymptotic independence implied by a Gaussian copula}
    \resumeItemListEnd

\resumeSubHeadingListEnd

%-----------EDUCATION-----------
\section{Education}
\resumeSubHeadingListStart
  \resumeSubheading
    {Master M2MO --- Quantitative Finance and Data Science}{Sep 2025 -- Present}
    {Universit\'e Paris Cit\'e}{Paris, France}
  \vspace{3pt}
  \resumeSubheading
    {MSc in Data Science and Risk Management}{Aug 2023 -- Dec 2025}
    {ENSAI}{Rennes, France}
  \vspace{3pt}
  \resumeSubheading
    {MSc in Statistics applied to Economics}{Sep 2021 -- Aug 2023}
    {ISSEA -- CEMAC}{Yaound\'e, Cameroon}
  \vspace{3pt}
  \resumeSubheading
    {Bachelor's Degree in Mathematics}{Sep 2017 -- Aug 2020}
    {Higher Teacher Training College \& University of Yaound\'e I}{Yaound\'e, Cameroon}
\resumeSubHeadingListEnd

%-----------EXPERIENCE-----------
\section{Research \& Professional Experience}
\resumeSubHeadingListStart
  \resumeSubheading
    {\textbf{Quantitative Analyst (Intern)} $|$ \textbf{PwC}}{Apr -- Oct 2026}
    {Risk Line of Service}{Paris, France}
    \resumeItemListStart
      \resumeItem{Default Risk Charge modelling under FRTB; research work described above}
    \resumeItemListEnd

  \resumeSubheading
    {\textbf{Quantitative Risk Analyst (Intern)} $|$ \textbf{Cr\'edit Agricole}}{Apr -- Sep 2025}
    {Risk Modeling Department}{Paris, France}
    \resumeItemListStart
      \resumeItem{Climate stress-testing methodologies using time series and Monte Carlo methods}
      \resumeItem{Automated backtesting tool in Python (Dash) for model validation}
    \resumeItemListEnd

  \resumeSubheading
    {\textbf{Intern Mathematics Teacher} $|$ \textbf{Ngoa Ekel\'e High School}}{Jan -- May 2022}
    {Mathematics Department}{Yaound\'e, Cameroon}
    \resumeItemListStart
      \resumeItem{Competency-based mathematics teaching; assessment design and performance analysis}
    \resumeItemListEnd
\resumeSubHeadingListEnd

%-----------SKILLS-----------
\section{Technical Skills}
\begin{itemize}[leftmargin=0.15in, label={}]
  \small{\item{
    \textbf{Languages:} Python (NumPy, SciPy, PyTorch, scikit-learn, statsmodels), R, C++ (basics), SQL, SAS, \LaTeX \\
    \textbf{Methods:} Stochastic control, HJB/HJBI PDEs, optimal stopping, particle filtering, Kalman filtering, Monte Carlo, spectral projected gradient, EM, random matrix theory, copulas, extreme value theory
  }}
\end{itemize}

%-----------REFERENCES-----------
\section{References}
\small
\textbf{Adrien Saumard} --- ENSAI, Rennes \hspace{1.2em}
\textbf{Basile De Loynes} --- ENSAI, Rennes \hspace{1.2em}
\textbf{Sim\'eon Fotso} --- \'Ecole Normale Sup\'erieure, Yaound\'e

\end{document}
